\begingroup
\setlength{\abovedisplayskip}{3pt plus 1pt minus 1pt}
\setlength{\belowdisplayskip}{3pt plus 1pt minus 1pt}
\setlength{\abovedisplayshortskip}{0pt plus 1pt}
\setlength{\belowdisplayshortskip}{3pt plus 1pt minus 1pt}
\pdfbookmark[2]{On IX.}{nexus-remarks-ix}
\begin{center}\large On IX.\end{center}

Let $p$ be a positive odd number divisible by no square, and put
\[
S_r=\Sigma\left(\frac{s_r}p\right)
\]
where $s_r$ must run through all numbers relatively prime to $p$ that lie between $(r-1)\frac p8$ and $r\frac p8$. Denote the numbers of properly primitive forms that are not properly equivalent, for the determinants $-p$ and $-2p$, by $C_1$ and $C_2$, respectively. Then (compare \textsc{Dirichlet}, Recherches sur diverses applications etc., § 11, in \textsc{Crelle}'s Journal XXI)
\[
C_1=2(S_1+S_2),\quad C_2=2(S_1-S_4)
\]
or
\[
C_1=S_1+S_2+S_3+S_4,\quad C_2=2(S_2+S_3)
\]
according as $p\equiv1$ or $\equiv3\pmod4$. Further, the numbers $s_1$ and $s_2$, taken together, form the same collection as the numbers $2s_1$ and $p-2s_4$; likewise the numbers $s_3$ and $s_4$ form the same collection as the numbers $2s_2$ and $p-2s_3$. In the case $p\equiv1\pmod4$, the sum $S_1+S_2+S_3+S_4=0$. These observations give two further relations between the four sums $S_1,S_2,S_3,S_4$ in both cases, so that each can be expressed in terms of $C_1$ and $C_2$. In this way, when $p\equiv1\pmod4$, one obtains
\[
\begin{aligned}
S_1=S_8&=\tfrac14\left(\frac2p\right)C_1+\tfrac14C_2\\
S_2=S_7&=\tfrac14\left(2-\left(\frac2p\right)\right)C_1-\tfrac14C_2\\
S_3=S_6&=-\tfrac14\left(2+\left(\frac2p\right)\right)C_1+\tfrac14C_2\\
S_4=S_5&=\tfrac14\left(\frac2p\right)C_1-\tfrac14C_2
\end{aligned}
\]
\SourceLeaf{312}
and, when $p\equiv3\pmod4$,
\[
\begin{aligned}
S_1=-S_8&=\tfrac14\left(3+\left(\frac2p\right)\right)C_1-\tfrac14C_2\\
S_2=-S_7&=-\tfrac14\left(1-\left(\frac2p\right)\right)C_1+\tfrac14C_2\\
S_3=-S_6&=\tfrac14\left(1-\left(\frac2p\right)\right)C_1+\tfrac14C_2\\
S_4=-S_5&=\tfrac14\left(1-\left(\frac2p\right)\right)C_1-\tfrac14C_2
\end{aligned}
\]

If $p$ is a prime, the formulas given in the text for the numbers of quadratic residues in the individual eighths follow immediately.
\BookRule{35mm}
\par\endgroup
